%% titre: Méthodes — Existence d'un minimum sur un ensemble non borné
%% matiere: Optimisation
%% niveau: Licence 3
%% annee: 2026-2027
%% ordre: 9
%% description: Garantir l'existence d'un minimum global sur un ensemble fermé non vide, même non borné, grâce à la coercivité.

\section{Théorème d'existence}

\begin{theorem}[Existence d'un minimum global]
Soit $S\subset\R^n$ un ensemble fermé et non vide, éventuellement non borné. Soit $f:S\to\R$ une fonction continue et coercive sur $S$, c'est-à-dire
\[
\lim_{\substack{\|x\|\to+\infty\\x\in S}}f(x)=+\infty.
\]
Alors il existe $x^*\in S$ tel que
\[
\boxed{f(x^*)=\min_{x\in S}f(x).}
\]
\end{theorem}

\section{Démarche}

\begin{methode}[Montrer l'existence d'un minimum global sur un ensemble non borné]
\begin{enumerate}[label=\textbf{Étape \arabic*.}]
    \item \textbf{Admissibilité} : vérifier que $S$ est non vide.

    \item \textbf{Fermeture} : montrer que $S$ est fermé.

    \item \textbf{Continuité} : vérifier que la fonction objectif $f$ est continue sur $S$.

    \item \textbf{Coercivité} : montrer que
    \[
    \lim_{\substack{\|x\|\to+\infty\\x\in S}}f(x)=+\infty.
    \]

    \item \textbf{Conclure} : d'après le théorème d'existence, il existe au moins un point $x^*\in S$ qui réalise le minimum global de $f$ sur $S$.
\end{enumerate}
\end{methode}

\section{Démonstration du théorème}

\begin{proof}
Puisque $S$ est non vide, choisissons $x_0\in S$.

Posons
\[
\alpha=f(x_0).
\]

La coercivité de $f$ garantit l'existence d'un réel $R>0$ suffisamment grand tel que
\[
x\in S,\quad \|x\|>R
\quad\Longrightarrow\quad
f(x)>\alpha.
\]

Considérons alors
\[
K=S\cap\overline B(0,R),
\]
où
\[
\overline B(0,R)=\{x\in\R^n:\|x\|\leq R\}.
\]

Puisque $S$ est fermé et que la boule fermée $\overline B(0,R)$ est compacte dans $\R^n$, leur intersection $K$ est compacte.

De plus, $K$ est non vide car on peut choisir $R$ suffisamment grand pour que $x_0\in K$.

La fonction $f$ étant continue sur $K$, le théorème de Weierstrass garantit l'existence de $x^*\in K$ tel que
\[
f(x^*)=\min_{x\in K}f(x).
\]

Or, pour tout $x\in S\setminus K$, on a
\[
f(x)>\alpha=f(x_0).
\]

Comme $x_0\in K$ et que $x^*$ minimise $f$ sur $K$,
\[
f(x^*)\leq f(x_0)=\alpha.
\]

Par conséquent, pour tout $x\in S\setminus K$,
\[
f(x)>f(x_0)\geq f(x^*).
\]

Ainsi, $x^*$ minimise également $f$ sur tout l'ensemble $S$.

Finalement,
\[
\boxed{\exists x^*\in S,\quad
f(x^*)=\min_{x\in S}f(x).}
\]
\end{proof}

\section{Exemple d'application}

\begin{exemple}
Considérons le problème
\[
(P)\qquad
\begin{cases}
\displaystyle\min_{(x,y)\in\R^2}
& f(x,y)=x^2+y^2-2x-4y\\
\text{s.c.}
& x+y\geq 1.
\end{cases}
\]

L'ensemble admissible est
\[
S=\{(x,y)\in\R^2:x+y\geq1\}.
\]

\textbf{1. Ensemble admissible}

On a
\[
S=g^{-1}([0,+\infty[),
\qquad g(x,y)=x+y-1.
\]

La fonction $g$ est continue et $[0,+\infty[$ est fermé. Donc $S$ est fermé.

De plus, $S$ est non vide et non borné.

\textbf{2. Continuité}

La fonction $f$ est polynomiale, donc continue sur $\R^2$ et en particulier sur $S$.

\textbf{3. Coercivité}

En complétant les carrés,
\[
f(x,y)=(x-1)^2+(y-2)^2-5.
\]

Ainsi,
\[
f(x,y)\geq \frac12(x^2+y^2)-5.
\]

Par conséquent,
\[
\lim_{\|(x,y)\|\to+\infty}f(x,y)=+\infty.
\]

La fonction $f$ est donc coercive sur $\R^2$, et par restriction, sur $S$.

\textbf{4. Conclusion}

L'ensemble $S$ est fermé et non vide, et la fonction $f$ est continue et coercive sur $S$.

D'après le théorème d'existence, le problème $(P)$ admet au moins une solution optimale globale.
\end{exemple}

\section{Attention aux hypothèses}

\begin{attention}
\begin{itemize}
    \item La fermeture de $S$ est essentielle pour garantir que la limite d'une suite admissible reste admissible.
    \item La coercivité doit être vérifiée uniquement sur $S$ : il n'est pas nécessaire que $f$ soit coercive sur tout $\R^n$.
    \item La continuité seule ne garantit pas l'existence d'un minimum sur un ensemble non borné.
    \item La coercivité seule ne suffit pas si $S$ n'est pas fermé.
    \item Ce théorème garantit l'existence d'au moins un minimum global, mais pas nécessairement son unicité.
\end{itemize}
\end{attention}
